% CP-VI-0046
% target_chapter: VI/33 — $\mathcal O_X$-Modules and Quasicoherent Sheaves
% source_unit_id: IMP-DL-0A9B1F0A75-P003
% source: Downloads/theory-of-algebraic-geometry-9.tex L1068-L1413
% solution_provenance: SOURCE_EXPLICIT_SOLUTION

\subsection*{Problem statement}

Let \(X\) be a scheme and let
\[
f:\mathcal{F}\longrightarrow \mathcal{G}
\]
be a morphism of quasi-coherent sheaves of \(\mathcal{O}_X\)-modules.

\begin{enumerate}[label=(\alph*)]
	\item Show that
	\[
	\ker f,\qquad \operatorname{coker} f,\qquad \operatorname{im} f
	\]
	are quasi-coherent.
	
	\item Further, if \(X\) is Noetherian and \(\mathcal{F}\) and \(\mathcal{G}\) are coherent, show that
	\[
	\ker f,\qquad \operatorname{coker} f,\qquad \operatorname{im} f
	\]
	are coherent.
	
	\item Let
	\[
	g:X\longrightarrow Y
	\]
	be a morphism of schemes, and let \(\mathcal{F}\) be a quasi-coherent, respectively coherent, sheaf of \(\mathcal{O}_Y\)-modules. Show that
	\[
	g^*\mathcal{F}
	\]
	is quasi-coherent, respectively coherent, on \(X\).
	
	\item Show by example that if \(\mathcal{G}\) is a coherent sheaf on \(X\), then
	\[
	g_*\mathcal{G}
	\]
	need not be coherent on \(Y\).
\end{enumerate}

% ------------------------------------------------------------
\subsection{Kernels, cokernels, and images of quasi-coherent sheaves}

\textbf{Solution.}

Let
\[
f:\mathcal{F}\longrightarrow \mathcal{G}
\]
be a morphism of quasi-coherent \(\mathcal{O}_X\)-modules.

Take an affine open subset
\[
U=\operatorname{Spec} A\subseteq X.
\]

Because \(\mathcal{F}\) and \(\mathcal{G}\) are quasi-coherent, there exist \(A\)-modules \(M\) and \(N\) such that
\[
\mathcal{F}|_U\cong \widetilde{M},
\qquad
\mathcal{G}|_U\cong \widetilde{N}.
\]

Since the association
\[
M\longmapsto \widetilde{M}
\]
gives an equivalence between \(A\)-modules and quasi-coherent sheaves on \(\operatorname{Spec} A\), the restricted morphism
\[
f|_U:\mathcal{F}|_U\longrightarrow \mathcal{G}|_U
\]
corresponds to an \(A\)-module homomorphism
\[
\phi:M\longrightarrow N.
\]

For modules on an affine scheme, localization is exact. Therefore sheafification preserves kernels, cokernels, and images:
\[
\ker(\widetilde{\phi})
\cong
\widetilde{\ker\phi},
\]
\[
\operatorname{coker}(\widetilde{\phi})
\cong
\widetilde{\operatorname{coker}\phi},
\]
and
\[
\operatorname{im}(\widetilde{\phi})
\cong
\widetilde{\operatorname{im}\phi}.
\]

Hence
\[
(\ker f)|_U
\cong
\widetilde{\ker\phi},
\]
\[
(\operatorname{coker} f)|_U
\cong
\widetilde{\operatorname{coker}\phi},
\]
and
\[
(\operatorname{im} f)|_U
\cong
\widetilde{\operatorname{im}\phi}.
\]

Each of these is quasi-coherent on the affine open subset \(U\).

Since quasi-coherence is a local property and the affine open subsets cover \(X\), we conclude that
\[
\boxed{
	\ker f,\qquad \operatorname{coker} f,\qquad \operatorname{im} f
	\text{ are quasi-coherent sheaves on }X.
}
\]

% ------------------------------------------------------------
\subsection{The coherent case on a Noetherian scheme}

\textbf{Solution.}

Now suppose that \(X\) is Noetherian and that
\[
\mathcal{F},
\qquad
\mathcal{G}
\]
are coherent sheaves.

Take an affine open subset
\[
U=\operatorname{Spec} A\subseteq X.
\]

Since \(X\) is Noetherian, the ring \(A\) is Noetherian. Since \(\mathcal{F}\) and \(\mathcal{G}\) are coherent, there exist finitely generated \(A\)-modules \(M\) and \(N\) such that
\[
\mathcal{F}|_U\cong \widetilde{M},
\qquad
\mathcal{G}|_U\cong \widetilde{N}.
\]

The restricted morphism is induced by an \(A\)-module homomorphism
\[
\phi:M\longrightarrow N.
\]

Because \(A\) is Noetherian:

\begin{itemize}
	\item every submodule of the finitely generated module \(M\) is finitely generated;
	\item every submodule of the finitely generated module \(N\) is finitely generated;
	\item every quotient of a finitely generated module is finitely generated.
\end{itemize}

Therefore
\[
\ker\phi,
\qquad
\operatorname{im}\phi,
\qquad
\operatorname{coker}\phi
\]
are finitely generated \(A\)-modules.

Consequently,
\[
\widetilde{\ker\phi},
\qquad
\widetilde{\operatorname{im}\phi},
\qquad
\widetilde{\operatorname{coker}\phi}
\]
are coherent sheaves on \(U\).

Since coherence is local, we conclude that
\[
\boxed{
	\ker f,\qquad \operatorname{coker} f,\qquad \operatorname{im} f
	\text{ are coherent sheaves on }X.
}
\]

\begin{remark}
	The Noetherian hypothesis is essential for kernels and images. Over a non-Noetherian ring, a submodule of a finitely generated module need not be finitely generated.
\end{remark}

% ------------------------------------------------------------
\subsection{Pullbacks of quasi-coherent and coherent sheaves}

\textbf{Solution.}

Let
\[
g:X\longrightarrow Y
\]
be a morphism of schemes, and let \(\mathcal{F}\) be a quasi-coherent sheaf of \(\mathcal{O}_Y\)-modules.

Recall that
\[
g^*\mathcal{F}
=
g^{-1}\mathcal{F}
\otimes_{g^{-1}\mathcal{O}_Y}
\mathcal{O}_X.
\]

Take an affine open subset
\[
V=\operatorname{Spec} A\subseteq Y.
\]

Since \(\mathcal{F}\) is quasi-coherent, there exists an \(A\)-module \(M\) such that
\[
\mathcal{F}|_V\cong \widetilde{M}.
\]

Now let
\[
U=\operatorname{Spec} B\subseteq g^{-1}(V)
\]
be an affine open subset of \(X\). The restricted morphism
\[
g|_U:U\longrightarrow V
\]
corresponds to a ring homomorphism
\[
A\longrightarrow B.
\]

The pullback of \(\widetilde{M}\) to \(U\) is given by
\[
(g^*\mathcal{F})|_U
\cong
\widetilde{M\otimes_A B}.
\]

Thus, on every affine open subset \(U\subseteq X\), the sheaf \(g^*\mathcal{F}\) is associated to a module. Hence
\[
\boxed{
	\mathcal{F}\text{ quasi-coherent on }Y
	\quad\Longrightarrow\quad
	g^*\mathcal{F}\text{ quasi-coherent on }X.
}
\]

Now suppose that \(\mathcal{F}\) is coherent.

Locally on \(V=\operatorname{Spec} A\), the module \(M\) may be taken to be finitely presented. Thus there exists an exact sequence
\[
A^p\longrightarrow A^q\longrightarrow M\longrightarrow 0
\]
with \(p,q<\infty\).

Tensoring with \(B\) gives an exact sequence
\[
B^p\longrightarrow B^q\longrightarrow M\otimes_A B\longrightarrow 0.
\]

Therefore
\[
M\otimes_A B
\]
is finitely presented as a \(B\)-module. Hence
\[
\widetilde{M\otimes_A B}
\]
is coherent on \(U\).

Since this holds locally on \(X\), we conclude that
\[
\boxed{
	\mathcal{F}\text{ coherent on }Y
	\quad\Longrightarrow\quad
	g^*\mathcal{F}\text{ coherent on }X.
}
\]

% ------------------------------------------------------------
\subsection{A pushforward of a coherent sheaf that is not coherent}

\textbf{Solution.}

Consider the structure morphism
\[
p:\mathbf{A}^1_k\longrightarrow \operatorname{Spec} k.
\]

Let
\[
\mathcal{G}=\mathcal{O}_{\mathbf{A}^1_k}.
\]

Since \(\mathbf{A}^1_k\) is affine and Noetherian, its structure sheaf is coherent.

We now compute the pushforward. Since \(\operatorname{Spec} k\) consists of a single point,
\[
p_*\mathcal{G}
\]
corresponds to the \(k\)-module
\[
\Gamma(\mathbf{A}^1_k,\mathcal{O}_{\mathbf{A}^1_k}).
\]

But
\[
\Gamma(\mathbf{A}^1_k,\mathcal{O}_{\mathbf{A}^1_k})
=
k[t].
\]

As a vector space over \(k\), the ring \(k[t]\) has basis
\[
1,t,t^2,t^3,\dots
\]
and is therefore infinite-dimensional.

A coherent sheaf on \(\operatorname{Spec} k\) corresponds to a finitely generated \(k\)-module, equivalently, to a finite-dimensional \(k\)-vector space. Since \(k[t]\) is not finite-dimensional over \(k\), the sheaf
\[
p_*\mathcal{O}_{\mathbf{A}^1_k}
\]
is not coherent.

Thus
\[
\boxed{
	\mathcal{O}_{\mathbf{A}^1_k}\text{ is coherent on }\mathbf{A}^1_k,
	\qquad
	p_*\mathcal{O}_{\mathbf{A}^1_k}\text{ is not coherent on }\operatorname{Spec} k.
}
\]

In particular, the pushforward of a coherent sheaf need not be coherent.

\begin{remark}
	In this example, the pushforward is nevertheless quasi-coherent:
	\[
	p_*\mathcal{O}_{\mathbf{A}^1_k}\cong \widetilde{k[t]}.
	\]
	The failure is not quasi-coherence, but finite generation over \(k\).
\end{remark}

% ============================================================
